% DF-001, CF-007/012, P09 del inventario REV05.
% Residencia: APP toroidal y transporte de cursor; ejemplo tras el cociclo.
% Propietarios: PAQUETE_ARTICULO_I_INTEGRACION_ESTADO_CAMPOS_20260919/
% deltas/integracion/APPRouteWinding.lean (íntegro),
% lean/biblioteca/TPKTransport.lean:19--143, 155--229,
% lean/biblioteca/TPKFiniteCursor.lean:17--146.
% Estatuto: RESULTADO_RECUPERADO; desarrollo del ejemplo finito efectivo,
% sin identificar el factor observable con todo el estado enriquecido.

\section[Retorno espacial y enrollamiento del cursor]{Retorno espacial y conservación del enrollamiento en el calendario de veintisiete transiciones}
\label{rc27:sec:cursor}

Las dos coordenadas periódicas de APP proceden de dos coordenadas
enteras transportadas. La lectura residual identifica posiciones que
difieren en múltiplos de nueve; los cocientes enteros conservan las
vueltas de ese transporte. El calendario TPK proporciona un ejemplo
completo en el que la posición residual retorna y el cursor levantado
conserva un desplazamiento horizontal distinto de cero.

\subsection{Coordenadas levantadas y cociclo de frontera}

Para $z\in\mathbb Z$, sea
\[
\rho_9(z)=1+((z-1)\bmod9),\qquad
q_9(z)=\left\lfloor\frac{z-1}{9}\right\rfloor.
\]
La división euclídea, también para coordenadas negativas, da
\begin{equation}
z=\rho_9(z)+9q_9(z),\qquad1\le\rho_9(z)\le9.
\label{rc27:eq:division}
\end{equation}
El acarreo de transporte de un incremento $a\in\mathbb Z$ es
\[
c(z,a)=q_9(z+a)-q_9(z).
\]
De \eqref{rc27:eq:division} se deducen
\begin{align}
\rho_9(z+a)+9c(z,a)&=\rho_9(z)+a,\\
c(z,a+b)&=c(z,a)+c(z+a,b).
\label{rc27:eq:cociclo}
\end{align}
La segunda identidad sigue de cancelar $q_9(z+a)$ en el miembro
derecho. El segundo incremento se evalúa desde el extremo del primero;
esa dependencia conserva la información del cruce de frontera.

Sea $D=\{N,E,S,O\}$, con desplazamientos
\[
\delta(N)=(-1,0),\quad\delta(E)=(0,1),\quad
\delta(S)=(1,0),\quad\delta(O)=(0,-1).
\]
Para una palabra de direcciones $w=d_1\cdots d_n$, definamos
$\Delta(w)=\sum_j\delta(d_j)$ y, desde $o=(o_1,o_2)$,
\[
\operatorname{Ext}(o,w)=o+\Delta(w),\qquad
\nu(o,w)=\bigl(c(o_1,\Delta_1(w)),c(o_2,\Delta_2(w))\bigr).
\]
El vector $\nu$ registra ambos enrollamientos.

\begin{proposition}[Composición e inversión de las vueltas]
Para dos rutas concatenables $u,v$ se cumple
\[
\nu(o,uv)=\nu(o,u)+\nu(\operatorname{Ext}(o,u),v).
\]
Si $w^{-1}$ invierte el orden y reemplaza cada dirección por su
opuesta, entonces
\[
\operatorname{Ext}(\operatorname{Ext}(o,w),w^{-1})=o,
\qquad
\nu(\operatorname{Ext}(o,w),w^{-1})=-\nu(o,w).
\]
Para una ruta espacialmente cerrada en APP,
\[
\nu(o,w)=\Delta(w)/9\in\mathbb Z^2.
\]
\end{proposition}
\begin{proof}
La suma de desplazamientos es aditiva bajo concatenación. Aplicar
\eqref{rc27:eq:cociclo} en cada coordenada demuestra la primera
identidad. La inversión da $\Delta(w^{-1})=-\Delta(w)$ y recupera
el extremo inicial. Restar los cocientes en orden contrario cambia
el signo de $\nu$. Por último, el cierre espacial equivale a
$\rho_9(o_i+\Delta_i)=\rho_9(o_i)$ para ambas coordenadas; la primera
igualdad de \eqref{rc27:eq:cociclo} da entonces $9\nu_i=\Delta_i$.
\end{proof}

Los ciclos de nueve movimientos al sur y de nueve movimientos al este,
desde $(1,1)$, tienen vueltas $(1,0)$ y $(0,1)$, respectivamente.
Son dos ciclos independientes de la superficie periódica APP.

\subsection{Actualización de los dos cursores}

Un cursor observable levantado es $c=(x,y,d)\in\mathbb Z^2\times D$.
La fase $p$ pertenece a $\mathbb Z/108\mathbb Z$ y se representa
por $0,\ldots,107$. En la hoja aditiva se activa el desplazamiento
cuando $p\equiv0\pmod3$; en la multiplicativa, cuando
$p\equiv1\pmod3$. Definamos
\[
a_+(p)=\mathbf1_{\{p\equiv0\ (3)\}},\qquad
a_\times(p)=\mathbf1_{\{p\equiv1\ (3)\}},\qquad
\varepsilon(p)=\mathbf1_{\{p+1\equiv0\ (27)\}}.
\]
El operador observable actúa mediante
\begin{equation}
c'_s=\bigl(x_s+a_s(p)\delta_1(d_s),\,
           y_s+a_s(p)\delta_2(d_s),\,
           \operatorname{op}^{\varepsilon(p)}d_s\bigr),
\qquad p'=p+1\pmod{108},
\label{rc27:eq:actualizacion}
\end{equation}
donde $s\in\{+,\times\}$ y $\operatorname{op}$ intercambia
$N$ con $S$ y $E$ con $O$.
El desplazamiento utiliza la dirección anterior y, después, la
transición de semiciclo actualiza la orientación.

La inversa recupera primero $p=p'-1\pmod{108}$, restituye
$d_s=\operatorname{op}^{\varepsilon(p)}d'_s$ y resta
$a_s(p)\delta(d_s)$ de las coordenadas. Como
$\operatorname{op}^2=I$, ambas composiciones son la identidad.

\begin{theorem}[Retorno espacial efectivo en la primera semicorona]
\label{rc27:thm:ejemplo}
Partamos de
$c_+(0)=c_\times(0)=(1,1,E)$ y $p(0)=0$.
Después de veintisiete transiciones de
\eqref{rc27:eq:actualizacion}, se obtiene
\[
c_+(27)=c_\times(27)=(1,10,O),\qquad p(27)=27.
\]
Ambas posiciones residuales vuelven a $(1,1)$, mientras que el
cociente horizontal aumenta en una unidad y el vertical permanece
invariable.
\end{theorem}
\begin{proof}
Para $0\le n\le27$, el número de activaciones aditivas anteriores
a la transición $n$ es $\lfloor(n+2)/3\rfloor$; el número de
activaciones multiplicativas es $\lfloor(n+1)/3\rfloor$. Antes
de la transición de fase $26$ a $27$ no hay inversión de dirección.
Por inducción en $n$, las coordenadas son
\[
x_+(n)=x_\times(n)=1,\qquad
y_+(n)=1+\left\lfloor\frac{n+2}{3}\right\rfloor,\qquad
y_\times(n)=1+\left\lfloor\frac{n+1}{3}\right\rfloor.
\]
Las direcciones son $E$ para $0\le n<27$ y $O$ para $n=27$.
La descomposición en cada grupo de tres transiciones es
\[
\begin{array}{c|c|c|c}
n&y_+(n)&y_\times(n)&\text{dirección}\\\hline
3k&1+k&1+k&E\\
3k+1&2+k&1+k&E\\
3k+2&2+k&2+k&E
\end{array}
\quad(0\le k\le8),
\]
y el extremo final es $(y_+,y_\times,d)=(10,10,O)$.
Hay nueve activaciones en cada hoja. La transición final tiene fase
previa $26$, de clase dos módulo tres; por ello invierte las
direcciones y mantiene las coordenadas que ya alcanzaron diez.

Finalmente, $\rho_9(10)=1$, $q_9(10)=1$,
$\rho_9(1)=1$ y $q_9(1)=0$. La posición residual coincide con
la inicial, y el enrollamiento espacial es $(0,1)$.
\end{proof}

El retorno de este ejemplo es espacial. La dirección y la fase
registradas son $(O,27)$ al final y $(E,0)$ al inicio. El ejemplo
conserva expresamente los tres datos: posición residual, orientación y
fase; su igualdad o diferencia se establece para cada uno.

\subsection{Recuperación del cursor y compatibilidad a toda profundidad}

\begin{proposition}[Recuperación del cursor levantado]
La posición residual, la dirección y los dos cocientes enteros
determinan unívocamente el cursor $c=(x,y,d)$.
\end{proposition}
\begin{proof}
Las dos identidades \eqref{rc27:eq:division} recuperan $x$ e $y$
y la dirección se conserva en la proyección observable. Si esos
cinco datos coinciden para dos cursores, coinciden sus tres
componentes y, por tanto, los cursores son iguales.
\end{proof}

Sea $P$ la proyección de los dos cursores a sus posiciones módulo
nueve, conservando sus direcciones y la fase. La actualización
finita $T_{\rm fin}$ usa los mismos indicadores de activación y
orientación, con suma residual en las coordenadas. Si $T$ es
\eqref{rc27:eq:actualizacion}, entonces
\[
PT=T_{\rm fin}P,\qquad PT^n=T_{\rm fin}^{\,n}P
\quad(n\ge0).
\]
La primera igualdad resulta de la compatibilidad de la reducción
módulo nueve con la suma y de la conservación de dirección y fase;
la segunda se demuestra por inducción en $n$.

El archivo ordenado de estados observables puede prolongarse mediante
$(q,\mathcal A)\mapsto(Tq,\mathcal A\Vert q)$.
Después de $n$ transiciones conserva su prefijo inicial y su longitud
aumenta en $n$. Esta operación mantiene separados el estado corriente,
el enrollamiento y el registro de sus estados anteriores. En la
construcción enriquecida completa, el factor observable se compone
con las demás lecturas de cilindro, incidencia y memoria mediante sus
proyecciones declaradas.
